% Evans 11.3, printed pp.595--601; English PDF pp.612--618.
% Chinese start: printed p.421 / scan p.439.
\section{Systems of Two Conservation Laws}
\label{sec:11-3}
\subsection{Local Riemann coordinates}
In this section $m=2$. Evans calls $w_i$ an $i$th Riemann invariant
when $Dw_i\cdot r_i=0$: it is constant on the state-space curve $R_i$.
A useful coordinate must also have $Dw_i\ne0$.
\begin{lemma}[Existence of local invariant coordinates; editorial detail]
Near a strictly hyperbolic state there are smooth $w_1,w_2$ satisfying
$Dw_i\cdot r_i=0$ and $\det D(w_1,w_2)\ne0$.
\end{lemma}
\begin{expandedproof}[flow-box construction]
Choose a short curve $\Gamma(a)$ transverse to $r_i$ at the base
point. Let $X(s,a)$ solve $X_s=r_i(X)$, $X(0,a)=\Gamma(a)$.
At the base point the two columns $r_i,\Gamma'$ are independent,
so the inverse function theorem makes $(s,a)$ local coordinates.
Define $w_i(X(s,a))=a$. Then $Dw_i r_i=0$ and $Dw_i\ne0$.
In two dimensions the annihilator of $r_i$ is spanned by $l_j$,
$j\ne i$. Thus $Dw_1$ and $Dw_2$ are nonzero multiples of
independent left eigenvectors. After intersecting their neighborhoods,
the two functions together are coordinates.
\end{expandedproof}

\begin{theorem}[Diagonal form; source \S11.3.1, Theorem 1]
For a smooth solution, set $v_i=w_i(u)$ and express $u=z(v_1,v_2)$.
Then
\begin{equation}\label{eq:11-diagonal}
 (v_1)_t+\lambda_2(v)(v_1)_x=0,\qquad
 (v_2)_t+\lambda_1(v)(v_2)_x=0.
\end{equation}
Here $\lambda_i(v)$ means $\lambda_i(z(v))$.
\end{theorem}
\begin{expandedproof}[why the indices cross]
Because $Dw_i$ is parallel to $l_j$, $j\ne i$,
$Dw_i(-DF+\lambda_jI)u_x=0$. The chain rule gives the displayed
transport equation. Along a curve $X'=\lambda_j(u(X,t))$ its $v_i$
value is constant. Also $z_{v_j}$ is a nonzero multiple of $r_i$
for $i\ne j$: it lies in the kernel of $Dw_i$ and cannot be zero
since $Dz$ is invertible. Therefore genuine nonlinearity of family
$i$ is equivalent to $\partial_{v_j}\lambda_i\ne0$. The transport
speeds still depend on both invariants, so the diagonal equations
are coupled, not independent scalar equations.
\end{expandedproof}
\PDEAnnotation{两种“不变”曲线}{按 Evans 的编号，$w_1$ 在状态空间
沿 $R_1$ 不变，但 $w_1(u(x,t))$ 沿时空的 $\lambda_2$ 特征线不变。
这两个陈述分别使用右特征向量与另一族左特征向量，并不矛盾。}

For barotropic gas in conserved variables $(\rho,q)$,
\[
 F=(q,q^2/\rho+p(\rho)),\qquad \rho>0,\quad p'>0.
\]
In primitive variables $(\rho,v)$, $v=q/\rho$, the matrix is
$\left(\begin{smallmatrix}v&\rho\\p'/\rho&v\end{smallmatrix}\right)$,
with speeds $v\mp c$, $c=\sqrt{p'}$. Put
\begin{equation}\label{eq:11-gas-invariants}
 I(\rho)=\int_{\rho_*}^{\rho}\frac{c(s)}s\,ds,\qquad
 w_1=I(\rho)+v,\quad w_2=I(\rho)-v.
\end{equation}
The Jacobian determinant in $(\rho,v)$ is $-2c/\rho\ne0$.
Direct substitution of
$\rho_t+v\rho_x+\rho v_x=0$ and
$v_t+vv_x+(c^2/\rho)\rho_x=0$ proves
$D_+w_1=0$, $D_-w_2=0$, where $D_\pm=\partial_t+(v\pm c)\partial_x$.
These equations use only the barotropic mass--momentum model.
It is not necessary to impose the full Euler energy equation with
constant internal energy.
\PDEAnnotation{远离真空的条件}{当 $\rho\downarrow0$，$c/\rho$
可能发散，特征间隙也可能消失。局部坐标可逆不保证整个解的值域
始终留在同一张坐标图中；下面的爆破定理显式规定紧矩形值域。}

\subsection{Compression and a weighted Riccati equation}
We give a version of source \S11.3.2, Theorem 2 with its uniformity
assumptions explicit. Suppose the characteristic variables exist on
an open set containing the closed rectangle
$K=[A_1,B_1]\times[A_2,B_2]$. Assume the coefficients are $C^2$ there,
\begin{equation}\label{eq:11-blowup-assumptions}
 \lambda_2-\lambda_1\ge d>0,\qquad
 \partial_{v_1}\lambda_2\ge c_0>0,\qquad
 \partial_{v_2}\lambda_1\ge c_0>0 \quad\hbox{on }K.
\end{equation}
The signs can be obtained from genuine nonlinearity by suitable
orientations of the two invariant coordinates on a connected chart.
\begin{theorem}[Finite-time loss of smoothness]
Suppose the smooth initial invariants range in $K$, and a classical
solution remains in this chart with range in $K$ while it exists.
If either initial derivative $(v_1)_x$ or $(v_2)_x$ is negative
somewhere, it cannot remain a classical solution for all positive time.
More precisely a weighted derivative diverges to $-\infty$ in a
finite bounded time unless the assumed smooth regime ends sooner.
\end{theorem}
\PDEAnnotation{爆破的是梯度}{输运使 $v_i$ 保持初值范围，
而 $e^H(v_1)_x$ 满足负二次 ODE。
证明使用的是系数的一致正下界；只有逐点正性而没有紧值域控制，
不能据此断言积分必在有限时刻达到临界值。}
\begin{expandedproof}[removing the mixed derivative term]
Set $a=(v_1)_x$, $b=(v_2)_x$ and $D_2=\partial_t+\lambda_2\partial_x$.
Differentiating the first equation in \eqref{eq:11-diagonal} gives
\[
 D_2a=-(\partial_1\lambda_2)a^2-(\partial_2\lambda_2)ab,
 \qquad D_2v_2=(\lambda_2-\lambda_1)b.
\]
Choose $H$ with
\[
 H(v_1,v_2)=\int_{A_2}^{v_2}
 \frac{\partial_2\lambda_2(v_1,s)}{\lambda_2(v_1,s)-\lambda_1(v_1,s)}\,ds.
\]
Along a $\lambda_2$ characteristic, $v_1$ is constant, so
$D_2H=(\partial_2\lambda_2)b$. For $y=e^H a$ the mixed term cancels:
\begin{equation}\label{eq:11-riccati}
 \frac{dy}{dt}=-k(t)y^2,\qquad
 k=e^{-H}\partial_1\lambda_2\ge\kappa>0.
\end{equation}
The lower bound follows from compactness of $K$ and
\eqref{eq:11-blowup-assumptions}, not merely pointwise positivity.
If $y(0)<0$, uniqueness preserves its sign before blow-up, and
\[
 \frac1{y(t)}=\frac1{y(0)}+\int_0^t k(s)\,ds.
\]
The denominator reaches zero no later than
$T_*=-1/(\kappa y(0))$. Thus a globally finite derivative is impossible.
For $b$, use $D_1=\partial_t+\lambda_1\partial_x$ and a weight
$J$ satisfying $\partial_1J=\partial_1\lambda_1/(\lambda_1-\lambda_2)$;
then $D_1(e^Jb)=-e^{-J}\partial_2\lambda_1(e^Jb)^2$.

On a smooth solution confined to the chart, the characteristic speeds
are bounded and locally Lipschitz in space, so the characteristics
exist on every finite smooth time strip. Each invariant is transported
from its initial range; this explains invariance of the rectangle
provided the chart covers that entire rectangle. It does not supply
a missing global inverse chart. Smooth compactly supported perturbations
of a constant state have bounded invariant ranges. Any nonconstant
one of these invariants, with equal limits at $\pm\infty$, has a
negative derivative somewhere, so the criterion applies under the
stated chart and uniformity assumptions.
\end{expandedproof}
\PDEBodyDiagram{11.3}{Compression reaches a vertical slope}
{An exact simple-wave example is the scalar Burgers characteristic map
$x=\alpha+t\,g(\alpha)$ for $g(\alpha)=-\alpha$ on a local interval.
Its Jacobian is $1-t$, and $u_x=-1/(1-t)$ until $t=1$.
The second panel plots this derivative. For two-field systems, the
weight in \eqref{eq:11-riccati} removes the mixed interaction and
recovers the same Riccati mechanism; the drawn map is the scalar model.}
{ch11-compression}
